% BHU PYQ -- Manually transcribed & error-corrected
% Paper: BCF-116 : Business Statistics
% Exam: B.Com. (Hons.) FMM Semester I Examination, 2022-23

\documentclass[11pt,a4paper]{article}
\usepackage[a4paper,top=2.5cm,bottom=2.5cm,left=2.8cm,right=2.8cm,headheight=14pt]{geometry}
\usepackage{amsmath,amssymb}
\usepackage[T1]{fontenc}
\usepackage[utf8]{inputenc}
\usepackage{lmodern,microtype,fancyhdr,enumitem,hyperref,lastpage,parskip,booktabs}

\hypersetup{
    colorlinks=true,
    urlcolor=black,
    linkcolor=black,
    pdftitle={BCF-116: Business Statistics -- BHU B.Com. (Hons.) FMM Sem I 2022-23}
}

\pagestyle{fancy}
\fancyhf{}
\renewcommand{\headrulewidth}{0.4pt}
\renewcommand{\footrulewidth}{0.4pt}
\fancyhead[L]{\small   \textit{Banaras Hindu University}}
\fancyhead[R]{\small   \textit{BCF-116: Business Statistics}}
\fancyfoot[C]{\small Page  \thepage\ of \pageref{LastPage}}


\newlist{parts}{enumerate}{2}
\setlist[parts,1]{label=(\alph*),leftmargin=*,itemsep=5pt,topsep=3pt}
\setlist[parts,2]{label=(\roman*),leftmargin=*,itemsep=3pt,topsep=2pt}

\newcommand{\pts}[1]{\hfill   \textit{[#1]}}


\begin{document}


\begin{center}
  {\large\bfseries BANARAS HINDU UNIVERSITY}\\[2pt]
  {
\normalsize B.Com. (Hons.) FMM --- Semester I Examination, 2022-23}\\[6pt]
  \rule{0.8\linewidth}{0.5pt}\\[6pt]
  {\large\bfseries Paper No. BCF-116: Business Statistics}\\[4pt]
  {
\normalsize Time: 3 Hours\hfill Maximum Marks: 70}
\end{center}

\rule{\linewidth}{0.4pt}

\smallskip



\noindent   \textit{Instructions: Answer all questions. All questions carry equal marks (14 marks each).}

\bigskip




\noindent   \textbf{Question 1.}
``Statistics is not a science, it is a scientific method.'' Discuss critically this statement and explain the utility and limitations of Statistics.\pts{14}

\centerline{   \textbf{OR}}

Rearrange the following series with equal intervals and then prepare 'less than' and 'more than' cumulative frequency distribution:\pts{14}

\begin{center}
  \small
  
\begin{tabular}{cc}
     oprule
   \textbf{Class Interval} &   \textbf{Frequency} \\
    \midrule
    0--5   & 3   \\
    5--6   & 2   \\
    6--9   & 7   \\
    9--12  & 5   \\
    12--17 & 16  \\
    17--18 & 12  \\
    18--20 & 15  \\
    20--24 & 20  \\
    24--25 & 8   \\
    25--30 & 10  \\
    30--36 & 2   \\
    \midrule
   \textbf{Total} &   \textbf{100} \\
    ottomrule
  \end{tabular}
\end{center}

\medskip\rule{\linewidth}{0.2pt}




\noindent   \textbf{Question 2.}
Define Geometric Mean and discuss its merits and demerits. Give two practical situations where you will recommend its use.\pts{14}

\centerline{   \textbf{OR}}

Find out mode and median wages from the following data:\pts{14}

\begin{center}
  \small
  
\begin{tabular}{lccccc}
     oprule
   \textbf{Wages ( extcurrency)} & 0--10 & 10--20 & 20--30 & 30--40 & 40--50 \\
    \midrule
   \textbf{No. of Workers}      & 5     & 7      & 15     & 25     & 8      \\
    ottomrule
  \end{tabular}
\end{center}

\medskip\rule{\linewidth}{0.2pt}




\noindent   \textbf{Question 3.}
Define Standard Deviation. Why is standard deviation regarded as superior to other measures of dispersion? Discuss.\pts{14}

\centerline{   \textbf{OR}}

Goals scored by teams A and B in a football season were as follows:\pts{14}

\begin{center}
  \small
  
\begin{tabular}{ccc}
     oprule
   \textbf{No. of Goals Scored in a Match} & \multicolumn{2}{c}{   \textbf{No. of Matches}} \\
    \cmidrule{2-3}
                                            &   \textbf{Team A} &   \textbf{Team B} \\
    \midrule
    0                                       & 27              & 17              \\
    1                                       & 9               & 9               \\
    2                                       & 8               & 6               \\
    3                                       & 5               & 5               \\
    4                                       & 4               & 3               \\
    ottomrule
  \end{tabular}
\end{center}
By calculating the coefficient of variation in each, state which team may be considered more consistent.

\medskip\rule{\linewidth}{0.2pt}




\noindent   \textbf{Question 4.}
What do you mean by index number? What points should be taken into consideration in the construction of index numbers? Explain.\pts{14}

\centerline{   \textbf{OR}}

From the following average prices of the commodities (given in rupees per unit), find chain base index no. with 2018 as base year:\pts{14}

\begin{center}
  \small
  
\begin{tabular}{cccccc}
     oprule
   \textbf{Items} &   \textbf{2018} &   \textbf{2019} &   \textbf{2020} &   \textbf{2021} &   \textbf{2022} \\
    \midrule
    I              & 8             & 10            & 12            & 15            & 18            \\
    II             & 4             & 5             & 8             & 10            & 12            \\
    ottomrule
  \end{tabular}
\end{center}

\medskip\rule{\linewidth}{0.2pt}




\noindent   \textbf{Question 5.}
What is time series analysis? Point out its importance.\pts{14}

\centerline{   \textbf{OR}}

Fit a straight line trend by the method of least squares from the following data:\pts{14}

\begin{center}
  \small
  
\begin{tabular}{lccccccc}
     oprule
   \textbf{Year}                  & 2010 & 2011 & 2012 & 2013 & 2014 & 2015 & 2016 \\
    \midrule
   \textbf{Production ('000 Ql.)} & 80   & 90   & 92   & 83   & 94   & 99   & 92   \\
    ottomrule
  \end{tabular}
\end{center}

\begin{parts}
  \item Fit a straight line by the least squares method and calculate the trend values.
  \item Also estimate the production for 2018.
  \item Show the above series and the trend on graph paper.
\end{parts}

\vfill
\rule{\linewidth}{0.4pt}

\begin{center}\small
\href{https://bhu-exam.vercel.app/}{Click here to visit our website}\\[4pt]
\href{https://me-aryan.vercel.app/}{Click here to visit our contributor}\\[4pt]
\href{https://quashan-me.vercel.app/}{Click here to visit our contributor}
\end{center}

\end{document}
